% !PRINTED 730
\weberpage{730}
%<<CONT
Classenzahl als Factor vorkommt. Die Classenzahl ist aber, da
gewiss immer wenigstens eine Classe, die Hauptclasse, existirt,
von Null verschieden, und also kann auch $T$ nicht gleich Null
sein. Wir haben daher den Satz:

1. \gesperrt{Die Zahlen $\tau_0, \tau_1, \ldots, \tau_{\nu'-1}$ sind ein unabhängiges
System von Normaleinheiten des Körpers $H_m$.}

Ist $\mathfrak{E}(r)$ eine Normaleinheit, so ist, wenn wir
\weq{7}{L_\beta = \log|\mathfrak{E}(r_\beta)|}
setzen, wegen (3)
\weq{8}{L_\beta = -L_{\beta+\nu'}.}
Ist nun $\mathfrak{E}_0, \mathfrak{E}_1, \ldots, \mathfrak{E}_{\nu'-1}$ irgend ein System unabhängiger
% [source note: the separator after \mathfrak{E}_0 is printed as a period-shaped or damaged comma; the mathematically required list punctuation is normalized to a comma.]
Normaleinheiten, so setzen wir:
\weq{9}{L_{i,k} = \log|\mathfrak{E}_i(r_k)|, \quad L_{i,k+\nu'} = -L_{i,k},}
so dass der Regulator des Systems der $\mathfrak{E}_i$ den Ausdruck erhält:
\weq{10}{\pm\Sigma\pm L_{0,0}L_{1,1}\cdots L_{\nu'-1,\nu'-1} = L(\mathfrak{E}_0, \mathfrak{E}_1, \ldots, \mathfrak{E}_{\nu'-1}),}
und eine von Null verschiedene Zahl ist, die wir mit $T_0$ bezeichnen.

Bedeutet $\mathfrak{E}(r)$ irgend eine andere Normaleinheit, so können
wir die Unbekannten $\xi_0, \xi_1, \ldots, \xi_{\nu'-1}$ aus den $\nu'$ linearen Gleichungen
 bestimmen, die wir erhalten, wenn wir in
\weq{11}{L_\beta = \xi_0 L_{0,\beta} + \xi_1 L_{1,\beta} + \cdots + \xi_{\nu'-1}L_{\nu'-1,\beta}}
den Index $\beta$ von 0 bis $\nu'-1$ gehen lassen, und wegen (8) und (9)
ist diese Gleichung auch für die übrigen Werthe von $\beta$ richtig,
so dass man daraus die für alle $r$ gültige Formel ableiten kann:
\weq{12}{\mathfrak{E}(r) = \pm\mathfrak{E}_0(r)^{\xi_0}\mathfrak{E}_1(r)^{\xi_1}\cdots\mathfrak{E}_{\nu'-1}(r)^{\xi_{\nu'-1}}.}
Bedeuten andererseits $m_0, m_1, \ldots, m_{\nu'-1}$ irgend welche ganze
rationale Zahlen, so sind alle in der Form
\[
\pm\mathfrak{E}_0(r)^{m_0}\mathfrak{E}_1(r)^{m_1}\cdots\mathfrak{E}_{\nu'-1}(r)^{m_{\nu'-1}}
\]
enthaltenen Zahlen Normaleinheiten des Körpers $H_m$. Demnach
lassen sich auf das Gleichungssystem (10) die Betrachtungen
anwenden, die wir im \S.\,186 ausführlich durchgeführt haben,
woraus sich ergiebt:

2. \gesperrt{Die aus (10) oder (11) bestimmten Exponenten
$\xi_0, \xi_1, \ldots, \xi_{\nu'-1}$ sind rationale Zahlen.}
